% Generated draft: requires mathematical and source review.
\documentclass[11pt]{article}
\usepackage[margin=25mm]{geometry}
\usepackage{fontspec}
\setmainfont{Latin Modern Roman}
\usepackage{amsmath,amssymb,mathtools,unicode-math}
\setmathfont{Latin Modern Math}
\usepackage{enumitem,needspace,xurl,hyperref}
\setlength{\parindent}{0pt}
\setlength{\parskip}{5pt}
\setlength{\emergencystretch}{3em}
\newcommand{\N}{\mathbb{N}}
\newcommand{\Z}{\mathbb{Z}}
\newcommand{\Q}{\mathbb{Q}}
\newcommand{\R}{\mathbb{R}}
\newcommand{\C}{\mathbb{C}}
\newcommand{\F}{\mathbb{F}}
\newcommand{\A}{\mathbb{A}}
\newcommand{\PP}{\mathbb{P}}
\newcommand{\E}{\mathbb{E}}
\newcommand{\eps}{\varepsilon}
\newcommand{\dd}{\,\mathrm{d}}
\begin{document}
\section*{UnsolvedMath ID: 20002712}
\textbf{Source problem code:} catalogue entry
\subsection{Definitions and mathematical statement}
Fix a finite Coxeter system \((W,S)\) and a reduced word \(c=s_1\cdots s_r\) for a Coxeter element. The \(c\)-sorting word of \(w\) is its lexicographically first reduced subword in \(c^\infty=c|c|\cdots\). The element is \(c\)-sortable if its sets of letters in successive blocks decrease by inclusion. Let \(T\) be the reflections and define absolute order by \(u\le_Tv\) when \(\ell_T(u)+\ell_T(u^{-1}v)=\ell_T(v)\); then \(NC(W,c)=[e,c]_T\).

For sorting word \(a_1\cdots a_\ell\), its inversion roots are \(\beta_i=a_1\cdots a_{i-1}(\alpha_{a_i})\), with corresponding reflections \(t_i\). Reading's \(nc_c\) is the ordered product of the cover reflections among these \(t_i\), where cover reflections are \(w s w^{-1}\) for right descents \(s\) of \(w\). The cluster map \(cl_c\) takes the inversion root at the last occurrence of each simple letter and the negative simple root for each absent letter, with the corresponding \(c\)-cluster compatibility convention. The original reference [41] is not reproduced, so its precise root/cluster labeling must be checked against these conventions.

Nathan Reading's four requests are: give a type-uniform count of sortable elements by \(\operatorname{Cat}(W)\); prove \(nc_c\) uniformly well-defined into \(NC(W,c)\); prove \(nc_c\) and \(cl_c\) uniformly bijective onto their noncrossing and cluster targets; and determine what happens in infinite type. For irreducible finite \(W\), \(\operatorname{Cat}(W)=\prod_{i=1}^r(h+d_i)/d_i\), with Coxeter number \(h\) and invariant degrees \(d_i\); use products over components otherwise.

The aid reports later Reading--Speyer Cambrian/arbitrary-Coxeter theory giving uniform finite-type map results, and a Galashin--Lam--Trinh--Williams type-uniform Hecke-theoretic Catalan enumeration. No full citations or hypotheses are supplied. These finite-type reports do not identify unchanged infinite-type targets automatically.

The explicit infinite-type test is
\[
 W=I_2(\infty)=\langle s,t:s^2=t^2=e\rangle,\qquad c=st.
\]
The finite-length \(c\)-sortable elements are \(e\), \(t\), and all nonempty alternating words starting with \(s\). The aid asserts that \(nc_c\) is injective and
\[
 \operatorname{im}(nc_c)=\{e\}\cup\{c^js:j\ge-1\},\qquad
 [e,c]_T=\{e,c\}\cup\{c^js:j\in\Z\}.
\]
Thus its image omits the Coxeter endpoint and one tail of reflections. Identify positive roots with their reflection labels and negative simples with signed labels \(-s,-t\). The cluster image is exactly
\[
 \{\{-s,-t\},\ \{-s,t\},\ \{s,-t\}\}
 \ \cup\ \{\{c^js,c^{j+1}s\}:j\ge0\}.
\]
It gives one ray in the bi-infinite rank-two exchange graph. Reversing orientation to \(c^{-1}=ts\) gives
\[
 \operatorname{im}(nc_{c^{-1}})=\{e\}\cup\{c^js:j\le0\}.
\]
The two reflection images cover all reflections and overlap at \(s\) and \(t=c^{-1}s\); both images also contain \(e\). Neither contains its own Coxeter endpoint. The opposite cluster orientation supplies the other ray. This is an exact rank-two boundary example, not a classification of every infinite Coxeter group or a proof of a universal infinite-type bijection.
\subsection{Short English statement}
The aid reports uniform finite-type counting and bijection results, then gives an infinite-dihedral example where the noncrossing and cluster maps remain meaningful but cover only one-sided targets. Reversing the Coxeter orientation supplies the opposite reflection tail and cluster ray.
\subsection{Sources}
\begin{itemize}
\item Local registry, permanent UnsolvedMath ID \(20002712\), exact public code \texttt{catalogue entry}; canonical location \texttt{aim-probability-notes.json}, \texttt{notes[153]}.
\item N. Reading, Problem 3.2 under ``3. Reflection Groups,'' AIM \emph{Braid groups, clusters and free probability}; supplied reference \url{https://aimath.org/WWN/braidgroups/braidgroups.pdf}. The verification field reports recovery of the split header ``Problem 3. 2.'' No PDF was accessed here. The source's reference [41] lacks its bibliographic entry in this record.
\item Reading--Speyer, Cambrian and arbitrary-Coxeter-group theory; Galashin--Lam--Trinh--Williams, uniform Hecke-theoretic Catalan enumeration, as attributed by the aid. No full titles, dates, journals, or theorem numbers are supplied.
\item Supplied research attempt 1: \texttt{partial\_result}, run and registry status \texttt{partially\_solved}, rigor label \texttt{proved}, low novelty confidence; classification \texttt{PARTIAL-PROGRESS}, created and updated 14 August 2026. The aid is machine generated, not peer reviewed, and requires independent verification before citation.
\end{itemize}
\end{document}
